\scanpage{72}
\[
\widetilde{\mathfrak g}\times\mathfrak g\longrightarrow\mathcal O_S\tag*{(41)}
\]
donnant donc des isom de Modules :
\[
\widetilde{\mathfrak g}\simeq\overset{\deleted{\unsure{\sim}}}{\mathfrak g}^\vee,\quad\mathfrak g\simeq\widetilde{\mathfrak g}^{\vee}\tag*{(42)}
\]

La \deleted{\unsure{a sur}} forme quadratique \inserted{$Q(u)=$} $-\det u$ sur $M$ est lisse, et la forme bilinéaire associée est
\[
\left\{\begin{aligned}
\Phi_1(\unsure{uv|u,v})&=-\operatorname{Tr}u\operatorname{Tr}v+\operatorname{Tr}(uv)=-\operatorname{Tr}u\operatorname{Tr}v+\Phi(u,v)\\
\text{donc}\quad-\det(u\unsure{+}v)+\det u+\det v&=\Phi_1(u,v)\\
&=-\operatorname{Tr}u\operatorname{Tr}v+\operatorname{Tr}(uv)
\end{aligned}\right.\tag*{(43)}
\]
qui est donc également dualisante. Sur $\widetilde{\mathfrak g}$ elle coïncide avec $\varphi$, et \deleted{\unsure{elle a}} comme $\operatorname{Tr}(1)=2$, on a
\[
\Phi_1(1,u)=-\Phi(1,u)=-\operatorname{Tr}u
\]
donc l'ortho. de $\mathcal O_S.1$ pour $\varphi_1$ est encore $\widetilde{\mathfrak g}$. L'accouplement \deleted{\unsure{le}} \inserted{$\mathfrak g\times\widetilde{\mathfrak g}\to\mathcal O_S$} déduit de $\varphi_1$ par passage au quotient \deleted{\unsure{diffère par le signe de}} \inserted{est le $=$ que celui} (41) défini par $\Phi$ ; \deleted{\unsure{puis que}} \inserted{D'ailleurs} \deleted{\unsure{la la r}} $\Phi$ et $\Phi_1$ ont $=$ restriction à $\widetilde{\mathfrak g}$, soit $\varphi$
\[
\left\{\begin{aligned}
\varphi&=\Phi\mid\widetilde{\mathfrak g}\times\widetilde{\mathfrak g}
=\Phi_1\mid\widetilde{\mathfrak g}\times\widetilde{\mathfrak g}
:\widetilde{\mathfrak g}\times\widetilde{\mathfrak g}\to\mathcal O_S\\
\varphi(uv)&=\operatorname{Tr}uv
\end{aligned}\right.\tag*{(42)}
\]
% archive sheet number: 71
